\subsection{对应于分次素理想的准素子模}

命题 3. 设 \(\Delta\) 是无挠交换群，A 是 \(a\) 分次诺特环，类型为 \(\Delta\)，\(\mathfrak{p}\) 是 A 的分次理想，M 是 \(a\) 分次 A-模，类型为 \(\Delta\)，且不归结为 0。假设对于 \(a\) 的每个齐次元 \(\mathfrak{p}\)，M 上比为 \(a\) 的伸缩映射是殆幂零的；对于 \(b\) 的每个齐次元 \(\mathfrak{p}\)，M 上比为 \(b\) 的伸缩映射是单射。则 \(\mathfrak{p}\) 是素理想，且 M 的子模 \(\{0\}\) 是 \(\mathfrak{p}\)-准素的。

只需证明 \(\operatorname{Ass}(\mathbf{M}) = \{\mathfrak{p}\}\)（\(\S 2\)，第 1 号，命题 1）。设 \(\mathfrak{q}\) 是与 \(\mathbf{M}\) 相伴的素理想；它是分次理想，且是某个齐次元 \(x \neq 0\) 的湮灭子，属于 \(\mathbf{M}\)（第 1 号，命题 1）。对于 \(a\) 的每个齐次元 \(\mathfrak{q}\)，都有 \(ax = 0\)，所以比为 \(a\) 的伸缩映射作用在 \(\mathbf{M}\) 上不是单射，从而 \(a \in \mathfrak{p}\)。反之，设 \(b\) 是 \(\mathfrak{p}\) 的一个齐次元；存在整数 \(n > 0\)，使得 \(b^n x = 0\)，于是 \(b^n \in \operatorname{Ann}(x) = \mathfrak{q}\)；由于 \(\mathfrak{q}\) 是素理想，\(b \in \mathfrak{q}\)。由于 \(\mathfrak{p}\) 和 \(\mathfrak{q}\) 分别由各自的齐次元生成，\(\mathfrak{p} = \mathfrak{q}\)，从而 \(\operatorname{Ass}(\mathbf{M}) \subset \{\mathfrak{p}\}\)。由于 \(\mathbf{M} \neq \{0\}\)，\(\operatorname{Ass}(\mathbf{M}) \neq \emptyset\)（\(\S 1\)，第 1 号，命题 2 的推论 1），所以 \(\operatorname{Ass}(\mathbf{M}) = \{\mathfrak{p}\}\)。

命题 4. 设 \(\Delta\) 是无挠交换群，A 是 \(\Delta\) 型分次诺特环，M 是 \(\Delta\) 型分次 A-模。设 \(\mathfrak{p}\) 是 A 的素理想，N 是 M 的关于 M 为 \(\mathfrak{p}\)-准素的子模。

\begin{enumerate}
\def\labelenumi{(\roman{enumi})}
\item
  \(\mathfrak{p}'\) 中包含的 A 的最大分次理想 \(\mathfrak{p}\)（代数，第二章，第 11 节，第 3 节）是素理想。
\item
  \(\mathbf{N}'\) 中最大的分次子模 \(\mathbf{N}\) 关于 \(\mathfrak{p}'\) 是 \(\mathbf{M}\)-准素的。
\end{enumerate}

  我们知道（同上），\(p'\)（分别为 \(N'\)）的齐次元恰好是 p（分别为 N）的齐次元。设 a 是 p 的齐次元；若 x 是 M 的齐次元，则存在整数 n \textgreater{} 0，使得 \(a^{n}x \in N\)；由于 \(a^{n}x\) 是齐次的，\(a^{n}x \in N'\)；由于每个 \(y \in M\) 都是有限个齐次元的直和，我们得到存在整数 q \textgreater{} 0，使得 \(a^{q}y \in N'\)，所以在 \(M/N'\) 中比为 a 的伸缩映射是殆幂零的。

  现在考虑 \(b\) 的一个齐次元 \(A - p'\)；由于 \(b \notin p\) 是齐次的，\(b\)。设 \(x\) 是 \(M\) 中满足 \(bx \in N'\) 的元素，并令 \((x_i)_{i \in \Delta}\) 为 \(x\) 的齐次分量族。由于 \(N'\) 是分次的，\(bx_i \in N'\) 对所有 \(i\) 都成立，所以 \(bx_i \in N\)；又由于 \(b \notin p\)，我们得到 \(x_i \in N\)；由于 \(x_i\) 是齐次的，\(x_i \in N'\)，从而 \(x \in N'\)，且比为 \(b\) 的伸缩映射作用在 \(M/N'\) 上是单射。命题 4 随即由命题 3 应用于 \(p'\) 和 \(M/N'\) 得出。

